
\section{Square-constraint labellings over \texorpdfstring{$\mathbb F_p$}{Fp}}
\label{sec:labellings}

\paragraph{Part IV.}
\Cref{sec:labellings} and the five sections after it prove the local
operator estimates with which Part~V verifies the hypotheses (CP1)--(CP7) of
the comparison theorem (\cref{thm:comparison}). The statements of Part~IV
are made for arbitrary oriented graphs, spectral data or couplings, and their
hypotheses are inequalities that a finite check verifies at one prime and that
the admissible majorants of \cref{lem:majorant} extend to all larger primes;
an evaluation at a particular graph and threshold, such as
\cref{lem:trace-mixing}\ref{it:trace:72}, is stated as an instance.

\paragraph{What this section proves and how to read it.}
This section proves the counting estimates behind every local law.
A \emph{graph labelling} is a map from the vertices of a small
oriented graph $G$ to $\mathbf F_p$ sending each directed edge to a
nonzero-square difference (\emph{injective} is a separate adjective;
$\lambda_p(G)$ is uniform on the injective labellings). The main
estimate is: once neighbours of a new vertex carry distinct labels,
the number of choices for the new label is $p2^{-t}$ up to relative
error $O(t2^tp^{-1/2})$ for $t$ constraints (Weil's bound).
Iterating gives sizes, marginals and conditional laws; comparing two
graphs gives a ratio of four counts; at $p\equiv3\pmod4$ exact
identities for directed triples are proved. Inputs: Weil's bound
(\cref{sec:notation}). Outputs: the counts used for cut operators,
spectral bounds and all Part~V local laws. A first reading should
follow the definitions and \cref{lem:extension}; sequential counts
can be used as reference. Terminology for Part~IV: a \emph{cut
operator} $A_p$ represents $E[F(x)H(x')]$ as
$\langle F,A_pH\rangle$ in the half law (defined in
\cref{sec:cuts}); a \emph{spectral cap} is a certified bound
$-gp^{-1/2}I\preceq A_p\preceq(1+c^+/p)I$ for all $p\ge p_0$ in a
stated residue class, for a stated threshold $p_0$;
\emph{pinning} means conditioning on prescribed labels; the
\emph{soft kernel} is the kernel obtained by expanding every
square indicator as $(1+\chi)/2$ and dropping the condition that
labels be distinct; the resulting terms in which two labels that
must be distinct coincide are the \emph{coincidence terms}, and
the difference between the true kernel and the soft kernel is a
sum of such terms. Full definitions are at the sites
noted; this paragraph only allows the statements to be read in order.

\paragraph{Conventions.}
Throughout, $p$ is an odd prime, $q=1/p$ and $r=p^{-1/2}$; $\chi=\chi_p$ is the
quadratic character of $\Fp$ with $\chi(0)=0$, and $\Qsq$ is the set of nonzero
squares. $L^2(\Fp)$ carries the uniform probability, $\E_x$ is the average over
$x\in\Fp$, $H_0\subseteq L^2(\Fp)$ is the space of mean-zero functions and
$\Jproj$ is the projection onto the constants. A strict-square pair is a pair
$(x,x')$ with $x'-x\in\Qsq$; $\Om$ is the set of them.

\subsection{Oriented graphs, labellings, and count errors}\label{subsec:labellings:defs}

The local laws of Part~V are uniform distributions on assignments
of elements of $\Fp$ to the vertices of a small directed graph,
with prescribed differences required to be nonzero squares. This
subsection defines those assignments (called labellings) and the
error terms that measure how close their counts are to what
independent uniform labels would give.

\begin{definition}[Oriented graphs and strict-square labellings]
\label{def:graphs:labeling}
An \emph{oriented graph} $G=(W,E)$ is a finite directed graph without loops
and with at most one edge on each unordered pair of vertices. The degree
$\deg_G(w)$ counts incoming and outgoing edges, and for $W'\subseteq W$,
$G[W']$ is the induced oriented graph and $G-w=G[W\setminus\{w\}]$. A
\emph{strict-square labelling} of $G$ over $\Fp$ is a map $z\colon W\to\Fp$ with
$z_v-z_u\in\Qsq$ for every edge $(u,v)\in E$; it is \emph{injective} if
$z_u\ne z_w$ for $u\ne w$. Let $\operatorname{Inj}_p(G)$ be the set of injective
strict-square labellings and, when it is nonempty, let $\lambda_p(G)$ (also
written $\lambda_G$) be the uniform probability on it: the \emph{graph law} of
$G$ at $p$. For a block $\Blk$ (a finite set of odd primes),
$\lambda_{\Blk}(G)=\bigotimes_{p\in\Blk}\lambda_p(G)$ is a law on
$(\prod_{p\in\Blk}\Fp)^W$; the labels at different primes are independent. A
labelling of a set $A\subseteq W$ is \emph{feasible} if it is the restriction
of an element of $\operatorname{Inj}_p(G)$.
\end{definition}

A vertex is labelled by a field element. The extension count below counts the
possible labels of one new vertex after its neighbours have distinct fixed
labels. Each neighbour imposes a condition on one quadratic character and so
halves the count, up to a character sum of size $O(\sqrt p)$: the main term is
$p2^{-t}$ for $t$ constraints. For integers $t,\Delta\ge0$ and $n,z\ge0$ put
\begin{equation}\label{eq:graphs:et}
 h_t=\frac t2-1+2^{-t},\qquad e_t(p,n)=2^t\bigl(h_tr+nr^2\bigr),\qquad
 w_\Delta(p,n)=2^\Delta\Bigl(\frac\Delta2r+nr^2\Bigr),
\end{equation}
the \emph{coarse} count error and its uniform majorant, and the \emph{graded}
coefficients and count errors
\begin{equation}\label{eq:graphs:graded}
\begin{gathered}
 a_t=(t-3)2^{t-1}+2,\qquad b_t=2^{t-1}-1\quad(t\ge1),\qquad a_0=b_0=0,\\
 e^+_t=a_tr+b_tr^2,\qquad e^-_t(z)=e^+_t+z2^tr^2 .
\end{gathered}
\end{equation}
Thus $a_0=a_1=a_2=0$, $a_3=2$, $a_4=10$, $a_5=34$, and $b_2=1$, $b_4=7$. We use
the following elementary properties.
\begin{enumerate}[label=\textup{(E\arabic*)}]
\item\label{it:et:mono} $a_{t+1}-a_t=(t-1)2^{t-1}$ and $b_{t+1}-b_t=2^{t-1}$
 for $t\ge1$, $a_1-a_0=b_1-b_0=0$, and $h_{t+1}-h_t=\frac12(1-2^{-t})$; so
 $a_t$, $b_t$, $h_t$ are nondecreasing in $t$, and $h_t\le t/2$.
\item\label{it:et:compare} $a_t+b_t=2^th_t$. Hence
 $e^+_t\le e^-_t(z)\le e_t(p,z)$, with $e_t(p,z)-e^-_t(z)=b_t(r-r^2)\ge0$, and
 $e_t(p,n)\le w_\Delta(p,n)$ for $t\le\Delta$.
\item\label{it:et:poly} For fixed $t,z,n,\Delta$, each of $e^+_t$, $e^-_t(z)$,
 $e_t(p,n)$, $w_\Delta(p,n)$ and its quotient by $r$ is a polynomial in $r$ with
 nonnegative coefficients; so all of them are nondecreasing in $r$, that is,
 nonincreasing in $p$. All are nondecreasing in $t$ ($w_\Delta$ in $\Delta$), $z$
 and $n$.
\item\label{it:et:main} $2^{-t}(1-e^-_t(z))$ and $2^{-t}(1-e_t(p,z))$ are
 nonincreasing in $t$ and in $z$.
\end{enumerate}
Indeed $a_t+b_t=(t-2)2^{t-1}+1=2^th_t$, which gives \ref{it:et:compare} with
\ref{it:et:mono}; \ref{it:et:poly} is read off \eqref{eq:graphs:et} and
\eqref{eq:graphs:graded}; and for \ref{it:et:main},
$2^{-t}(1-e^-_t(z))=2^{-t}-2^{-t}a_tr-2^{-t}b_tr^2-zr^2$ and
$2^{-t}(1-e_t(p,z))=2^{-t}-h_tr-zr^2$, where $2^{-t}a_t$ and $2^{-t}b_t$ are
nondecreasing (for $t\ge1$ their increments are $\frac12-2^{-t}$ and
$2^{-t-1}$, and from $t=0$ to $t=1$ they are zero).

\subsection{Counting choices for one new vertex}\label{subsec:labellings:extension}

An edge between a new vertex, labelled $x$, and a vertex labelled $\beta$
imposes $x-\beta\in\Qsq$ (edge into the new vertex) or $\beta-x\in\Qsq$ (edge
out of it), that is, $\chi(x-\beta)=1$ or $\chi(x-\beta)=\chi(-1)$. The
admissible labels of a vertex whose $t$ neighbours carry distinct labels, and
which must avoid a set of used labels, therefore form a set $\Phi$ as in the
following lemma, with the neighbour labels as roots. More generally, for
$\epsilon_i\in\{\pm1\}$ the condition $\epsilon_i(x-\beta_i)\in\Qsq$ is
$\chi(x-\beta_i)=\chi(\epsilon_i)$, so counts of this form are the case
$s_i=\chi(\epsilon_i)$. At $p\equiv3\pmod4$ the sign records the orientation of
the edge; at $p\equiv1\pmod4$ every sign is $+1$.

\begin{lemma}[Graded extension count]\label{lem:extension}
Let $t\ge0$, let $\beta_1,\dots,\beta_t\in\Fp$ be distinct, let
$s_1,\dots,s_t\in\{1,-1\}$, and let $\mathcal X\subseteq\Fp$ contain
$\beta_1,\dots,\beta_t$. Put
\[
 \Phi=\{x\in\Fp\setminus\mathcal X:\ \chi(x-\beta_i)=s_i\ (1\le i\le t)\},\qquad
 \mathfrak S=\sum_{x\in\Fp}\prod_{i=1}^t\bigl(1+s_i\chi(x-\beta_i)\bigr).
\]
\begin{enumerate}[label=\textup{(\alph*)}]
\item\label{it:ext:soft} $|\mathfrak S-p|\le a_t\sqrt p+b_t$.
\item\label{it:ext:upper} $|\Phi|\le p2^{-t}(1+e^+_t)$, whatever the size of
 $\mathcal X$.
\item\label{it:ext:lower} If $|\mathcal X|\le z$, then $|\Phi|\ge p2^{-t}(1-e^-_t(z))$.
\item\label{it:ext:coarse} If $|\mathcal X|\le n$, then
 $p2^{-t}-h_t\sqrt p-n\le|\Phi|\le p2^{-t}+h_t\sqrt p$; hence
 $|\Phi|\in p2^{-t}[1-e_t(p,n),1+e_t(p,n)]$.
\end{enumerate}
\end{lemma}

\begin{proof}
For $y\ne0$ the indicator of $\chi(y)=s$ is $\frac12(1+s\chi(y))$. Let
$f(x)=2^{-t}\prod_i(1+s_i\chi(x-\beta_i))$, so that $\mathfrak S=2^t\sum_xf(x)$.
Each factor $\frac12(1+s_i\chi(x-\beta_i))$ lies in $\{0,\frac12,1\}$, so
$0\le f\le1$, and for $x\notin\{\beta_1,\dots,\beta_t\}$, $f(x)$ is the
indicator of the $t$ conditions. Since the roots lie in $\mathcal X$,
\begin{equation}\label{eq:ext:split}
 |\Phi|=\sum_{x\notin\mathcal X}f(x)=2^{-t}\mathfrak S-\sum_{x\in\mathcal X}f(x),
 \qquad 0\le\sum_{x\in\mathcal X}f(x)\le|\mathcal X| .
\end{equation}
Expanding the product and using the multiplicativity of $\chi$,
\[
 \mathfrak S=\sum_{S\subseteq\{1,\dots,t\}}s_SW_S,\qquad
 s_S=\prod_{i\in S}s_i,\qquad
 W_S=\sum_{x\in\Fp}\chi\Bigl(\prod_{i\in S}(x-\beta_i)\Bigr).
\]
Here $W_\emptyset=p$; $W_S=\sum_x\chi(x-\beta)=0$ for $|S|=1$; $W_S=-1$ for
$|S|=2$ by \eqref{eq:two-root}; for odd $|S|=k\ge3$, $|W_S|\le(k-1)\sqrt p$ by
\cref{thm:weil}\textup{(a)}, the roots being distinct; and for even $|S|=k\ge4$,
$|W_S+1|\le(k-2)\sqrt p$ by \cref{thm:weil}\textup{(b)}, so
$|W_S|\le(k-2)\sqrt p+1$. Summing,
\[
 |\mathfrak S-p|\le\sqrt p\Bigl[\sum_{k\ \mathrm{odd},\,k\ge3}\binom tk(k-1)
 +\sum_{k\ \mathrm{even},\,k\ge4}\binom tk(k-2)\Bigr]
 +\sum_{k\ \mathrm{even},\,k\ge2}\binom tk .
\]
For $t\ge1$ the last sum is $2^{t-1}-1=b_t$. The bracket equals
$\sum_{k\ge3}\binom tk(k-1)-\sum_{k\ \mathrm{even},\,k\ge4}\binom tk$; since
$\sum_{k\ge0}\binom tk(k-1)=(t-2)2^{t-1}$ and
$\sum_{k\ \mathrm{even}}\binom tk=2^{t-1}$, removing the terms $k\le2$ gives
$\bigl[(t-2)2^{t-1}+1-\binom t2\bigr]-\bigl[2^{t-1}-1-\binom t2\bigr]=a_t$. This
proves \ref{it:ext:soft}; for $t=0$, $\mathfrak S=p$.

By \ref{it:ext:soft}, $2^{-t}\mathfrak S\in p2^{-t}[1-e^+_t,1+e^+_t]$. The
upper bound of \eqref{eq:ext:split} uses only $f\ge0$, which gives
\ref{it:ext:upper}; its lower bound gives
$|\Phi|\ge p2^{-t}(1-e^+_t)-z=p2^{-t}(1-e^-_t(z))$, which is \ref{it:ext:lower}.
Finally $2^{-t}(a_t\sqrt p+b_t)\le2^{-t}(a_t+b_t)\sqrt p=h_t\sqrt p$ by
\ref{it:et:compare}, and \ref{it:ext:coarse} follows from \ref{it:ext:upper}
and \ref{it:ext:lower}.
\end{proof}

Two features of the count recur below. First, the excluded labels are handled
collectively, through \eqref{eq:ext:split}: the exact value $W_S=-1$ of a pair
sum is available only for the sum over all of $\Fp$, and the partial sums of
$\chi(x-\beta_i)\chi(x-\beta_j)$ over $x\notin\mathcal X$ are not $-1$ when
$\mathcal X$ contains more than the roots. Second, Weil's bound needs distinct
roots. This is why every count in this paper is a count of \emph{injective}
extensions: the pinned labels are then distinct, and injectivity contributes only the excluded-label term $z2^tr^2$ of $e^-_t(z)$, which is of lower order.
The coarse error $e_t(p,n)$ results from bounding $|W_S|$ for $|S|=2$ by $\sqrt p$ instead of using $W_S=-1$, and for even $|S|=k\ge4$ by $(k-1)\sqrt p$; it is larger than the graded error by
\ref{it:et:compare}, so every statement proved or certified with $e_t(p,n)$
remains valid, and the certificates use both.

\begin{corollary}[Signed neighbourhoods]\label{lem:record:signed-nbhd}
Let $x_1,\dots,x_4\in\Fp$ be distinct, $s\in\{\pm1\}^4$, and put
$\upsilon_+=11r+6q$ and $\upsilon_-=11r+70q$. The soft count
$\mathfrak S=\sum_{z\in\Fp}\prod_i(1+s_i\chi(z-x_i))$ satisfies
$|\mathfrak S-p|\le10\sqrt p+7\le11\sqrt p+6$, and the signed neighbourhood
$\{z\notin\{x_1,\dots,x_4\}:\chi(z-x_i)=s_i\ \forall i\}$ has between
$\frac p{16}(1-\upsilon_-)$ and $\frac p{16}(1+\upsilon_+)$ elements. If $k$
further labels are excluded, at least $\frac p{16}(1-\upsilon_--16kq)$ of its
elements remain.
\end{corollary}

\begin{proof}
Take $t=4$ in \cref{lem:extension}, so $a_4=10$, $b_4=7$, and $\mathcal X=\{x_1,\dots,x_4\}$.
Part \ref{it:ext:soft} gives $|\mathfrak S-p|\le10\sqrt p+7\le11\sqrt p+6$, part
\ref{it:ext:upper} gives at most $\frac p{16}(1+10r+7q)\le\frac p{16}(1+\upsilon_+)$
elements and part \ref{it:ext:lower} with $z=4$ at least
$\frac p{16}(1-10r-71q)\ge\frac p{16}(1-\upsilon_-)$, since $q\le r$. Removing $k$
labels costs at most $k=\frac p{16}\cdot16kq$.
\end{proof}

\subsection{Sequential counts and densities}\label{subsec:labellings:sequential}

Labelling the vertices of a graph one at a time and applying
\cref{lem:extension} at each step bounds the number of completions of a
partial labelling. The lemmas of this subsection hold with any count error
that the extension count supplies; we isolate in \cref{def:graphs:error-pair} the property of the error that the lemmas use.

\begin{definition}[One-step error pairs]\label{def:graphs:error-pair}
Let $n\ge1$. A \emph{one-step error pair} at $(p,n)$ consists of two sequences
$(\epsilon^+_t)_{t\ge0}$, $(\epsilon^-_t)_{t\ge0}$ of nonnegative numbers,
nondecreasing in $t$, such that
$p2^{-t}(1-\epsilon^-_t)\le|\Phi|\le p2^{-t}(1+\epsilon^+_t)$ for every set $\Phi$ of
\cref{lem:extension} with $t$ roots and an excluded set of at most $n$ labels.
\end{definition}

By \cref{lem:extension} and \ref{it:et:mono}--\ref{it:et:poly}, two one-step
error pairs at $(p,n)$ are the \emph{graded pair}, $\epsilon^+_t=e^+_t$ and
$\epsilon^-_t=e^-_t(n)$, and the \emph{coarse pair}, $\epsilon^\pm_t=e_t(p,n)$;
so is every pair of larger nondecreasing sequences.

\begin{lemma}[Sequential counts]\label{lem:graphs:sequential}
Let $G=(W,E)$ be an oriented graph, $n\ge|W|$, $A\subseteq W$, and let
$v_1,\dots,v_s$ be an ordering of $W\setminus A$. Let $t_i$ be the number of
$G$-neighbours of $v_i$ in $A\cup\{v_1,\dots,v_{i-1}\}$, so that
$\sum_it_i=|E|-|E(G[A])|$, and put $\mathrm{main}=p^s2^{-\sum_it_i}$. Let
$(\epsilon^+,\epsilon^-)$ be a one-step error pair at $(p,n)$ with
$\epsilon^-_{t_i}<1$ for all $i$, and define the \emph{logarithmic count errors}
of the ordering,
\[
 \mathrm{up}=\sum_i\log\bigl(1+\epsilon^+_{t_i}\bigr),\qquad
 \mathrm{lo}=\sum_i-\log\bigl(1-\epsilon^-_{t_i}\bigr).
\]
For $x\in\operatorname{Inj}_p(G[A])$ let $C(x)$ be the number of
$y\in\operatorname{Inj}_p(G)$ with $y|_A=x$.
\begin{enumerate}[label=\textup{(\alph*)}]
\item\label{it:seq:count} $\mathrm{main}\cdot e^{-\mathrm{lo}}\le C(x)\le\mathrm{main}\cdot e^{\mathrm{up}}$
 for every $x\in\operatorname{Inj}_p(G[A])$; the upper bound holds without the
 hypothesis on $\epsilon^-$. In particular every element of
 $\operatorname{Inj}_p(G[A])$ is feasible, and $\operatorname{Inj}_p(G)\ne\emptyset$
 if $\operatorname{Inj}_p(G[A])\ne\emptyset$ (for instance if $A=\emptyset$).
\item\label{it:seq:marg} The marginal of $\lambda_p(G)$ on $A$ is supported on
 $\operatorname{Inj}_p(G[A])$, and its density with respect to $\lambda_p(G[A])$
 takes values in $[e^{-(\mathrm{up}+\mathrm{lo})},e^{\mathrm{up}+\mathrm{lo}}]$.
\item\label{it:seq:rows} \emph{Anchored extension probability.} Let $Z=W\setminus A$ and let $H$ be the
 set of edges of $G$ between $A$ and $Z$. For every $x\in\operatorname{Inj}_p(G[A])$,
 the probability under $y\sim\lambda_p(G[Z])$ that $(x,y)\in\operatorname{Inj}_p(G)$
 lies in $2^{-|H|}[e^{-(\mathrm{up}+\mathrm{lo})},e^{\mathrm{up}+\mathrm{lo}}]$.
\item\label{it:seq:coarse} For the coarse pair,
 $e^{\mathrm{up}+\mathrm{lo}}=R:=\prod_i\frac{1+e_{t_i}(p,n)}{1-e_{t_i}(p,n)}$. If
 $G$ has maximum degree $\Delta$ and $e_\Delta(p,n)<1$, the hypotheses hold for
 the coarse pair, for every $A$ and every ordering.
\item\label{it:seq:thresholds} \emph{Thresholds.} For the graded and the coarse
 pairs, the quantities $\epsilon^\pm_{t_i}$, $\mathfrak u=\sum_i\epsilon^+_{t_i}/r$
 and $\mathfrak l=\mathrm{lo}/r$, and for the coarse pair also $\log R/r$, are
 nondecreasing functions of $r$ on the range where every $\epsilon^-_{t_i}<1$.
 Hence, if $\epsilon^-_{t_i}<1$ at a real threshold $P>1$, then the same holds at
 every prime $p\ge P$, and there
 $\log(\max_xC(x)/\mathrm{main})\le r_p\,\mathfrak u(P)$ and
 $\log(\mathrm{main}/\min_xC(x))\le r_p\,\mathfrak l(P)$, where $r_p=p^{-1/2}$ and
 $\mathfrak u(P)$, $\mathfrak l(P)$ are evaluated at $r=P^{-1/2}$.
\end{enumerate}
\end{lemma}

\begin{proof}
Every edge not inside $A$ is counted exactly once in $\sum_it_i$, at its
endpoint of largest index; this gives $\sum_it_i=|E|-|E(G[A])|$.

\ref{it:seq:count} The elements of $\operatorname{Inj}_p(G)$ extending $x$ are
obtained by choosing $y(v_1),y(v_2),\dots$ in turn. When $v_i$ is labelled,
its $t_i$ labelled neighbours carry pairwise distinct labels (injectivity),
and, since $G$ has at most one edge on each unordered pair, each of them
imposes exactly one condition $\chi(y(v_i)-\beta)=s$, with $\beta$ its label. The
new label must avoid the $|A|+i-1\le n-1$ labels used so far, which contain
the roots. Conversely every sequence of choices satisfying these conditions
gives an element of $\operatorname{Inj}_p(G)$ extending $x$: the edges inside
$A$ are satisfied by $x$, every other edge at the step of its later endpoint.
So the number of choices at step $i$ is the size of a set $\Phi$ of
\cref{lem:extension} with $t_i$ roots and at most $n-1$ excluded labels, and it
lies in $p2^{-t_i}[1-\epsilon^-_{t_i},1+\epsilon^+_{t_i}]$ whatever the earlier
choices were. Multiplying over $i$ gives \ref{it:seq:count}.

\ref{it:seq:marg} The marginal of $\lambda_p(G)$ at $x$ is
$C(x)/\sum_{x'}C(x')$ and $\lambda_p(G[A])(x)=1/|\operatorname{Inj}_p(G[A])|$, so
the density is $C(x)/\operatorname{mean}_{x'}C(x')$, which lies between
$\min C/\max C$ and $\max C/\min C$, hence in the stated interval by
\ref{it:seq:count}; it is positive everywhere on $\operatorname{Inj}_p(G[A])$.

\ref{it:seq:rows} The probability is $C(x)/|\operatorname{Inj}_p(G[Z])|$. By
\ref{it:seq:count}, $C(x)\in p^s2^{-|H|-|E(G[Z])|}[e^{-\mathrm{lo}},e^{\mathrm{up}}]$.
By \ref{it:seq:count} for $G[Z]$ with no anchors and the same ordering,
$|\operatorname{Inj}_p(G[Z])|\in p^s2^{-|E(G[Z])|}\prod_i[1-\epsilon^-_{t^0_i},1+\epsilon^+_{t^0_i}]$,
where $t^0_i\le t_i$ is the number of earlier neighbours of $v_i$ in $Z$; since
$\epsilon^\pm$ are nondecreasing in $t$, this interval lies in
$p^s2^{-|E(G[Z])|}[e^{-\mathrm{lo}},e^{\mathrm{up}}]$. Divide.

\ref{it:seq:coarse} The identity is the definition. If $e_\Delta(p,n)<1$ then
$e_{t_i}(p,n)\le e_\Delta(p,n)<1$ by \ref{it:et:poly}, since $t_i\le\Delta$.

\ref{it:seq:thresholds} By \ref{it:et:poly}, $\epsilon^\pm_t$ and $\epsilon^\pm_t/r$ are
polynomials in $r$ with nonnegative coefficients. The functions
$x\mapsto-\log(1-x)/x=\sum_{k\ge1}x^{k-1}/k$ and
$x\mapsto\operatorname{artanh}(x)/x=\sum_{k\ge0}x^{2k}/(2k+1)$ are increasing
on $(0,1)$, so
$-\log(1-\epsilon^-)/r=\bigl(-\log(1-\epsilon^-)/\epsilon^-\bigr)(\epsilon^-/r)$ and
$\log\frac{1+e}{1-e}/r=2\bigl(\operatorname{artanh}(e)/e\bigr)(e/r)$ are products of
nonnegative nondecreasing functions of $r$ (this is an instance of
\cref{lem:majorant}). Since $r_p\le P^{-1/2}$ for $p\ge P$, the hypothesis
transfers from $P$ to $p$, and by \ref{it:seq:count} and $\log(1+x)\le x$,
$\log(\max C/\mathrm{main})\le\sum_i\epsilon^+_{t_i}(p)=r_p\mathfrak u(p)\le
r_p\mathfrak u(P)$; the other bound is the same argument.
\end{proof}

The next lemma collects the facts about one vertex, about subsets and about
conditional laws of a graph law that the cut and spectral estimates use; its
part \ref{it:vertex:tuple} is the density bound behind the extraction of
\cref{lem:matching:extractionk}.

\begin{lemma}[Vertices, subsets and conditional laws]\label{lem:graphs:vertex}
Let $G=(W,E)$ be an oriented graph with $|W|\le n$, maximum degree $\Delta$ and
$e_\Delta(p,n)<1$. Let $\lambda=\lambda_p(G)$, $d_w=\deg_G(w)$, and for $Y\subseteq W$
put
\[
 R_Y=\prod_{i\in Y}\frac{1+e_{d_i}(p,n)}{1-e_{d_i}(p,n)},\qquad d_i=\deg_G(i).
\]
\begin{enumerate}[label=\textup{(\alph*)}]
\item\label{it:vertex:inv} $\lambda$ is invariant under the maps
 $z\mapsto(\alpha z_w+\beta)_{w\in W}$ with $\alpha\in\Qsq$, $\beta\in\Fp$, and its
 one-vertex marginals are uniform (\cref{lem:marginals}).
\item\label{it:vertex:dens} Let $w\in W$. Conditionally on $z_{W\setminus\{w\}}$, the
 label $z_w$ is uniform on its admissible set $\Phi_w(z_{W\setminus\{w\}})$, which
 has at least $p2^{-d_w}(1-e_{d_w}(p,n))$ elements; so its conditional density
 with respect to the uniform probability on $\Fp$ is at most
  $2^{d_w}/(1-e_{d_w}(p,n))$. If $m$ is a probability measure on
  $\operatorname{Inj}_p(G)$ whose density with respect to $\lambda$ takes values in
  an interval $[c,Rc]$, $c>0$, then under $m$ the conditional probability of each
 value of $z_w$ given $z_{W\setminus\{w\}}$ is at most $R/|\Phi_w|$.
\item\label{it:vertex:del} The marginal of $\lambda$ on $W\setminus\{w\}$ is
 supported on $\operatorname{Inj}_p(G-w)$, and its density with respect to
 $\lambda_p(G-w)$ takes values in $[R_{\{w\}}^{-1},R_{\{w\}}]$.
\item\label{it:vertex:subset} Let $J\subseteq W$. The marginal $\lambda_J$ of
 $\lambda$ on $J$ is supported on $\operatorname{Inj}_p(G[J])$, and its density with
 respect to $\lambda_p(G[J])$ takes values in an interval $[c',R_{W\setminus J}c']$,
 $c'>0$ (in particular in $[R_{W\setminus J}^{-1},R_{W\setminus J}]$). For $m$ as in
 \ref{it:vertex:dens}, the density of the marginal $m_J$ with respect to
 $\lambda_p(G[J])$ takes values in an interval $[c'',R\,R_{W\setminus J}c'']$.
\item\label{it:vertex:tuple} Let $W=V\sqcup C\sqcup U$ with $V\ne\emptyset$, let $x_C$ be a
 feasible labelling of $C$, order $V$, and let $t_i$ be the number of
 neighbours of the $i$-th vertex of $V$ in $C$ and among the earlier vertices of
 $V$, so that $t_V=\sum_it_i$ is the number of edges inside $V$ or between $V$
 and $C$. The conditional law of $z_V$ given $z_C=x_C$ is supported on injective
 tuples, and its density with respect to the uniform probability on $\Fp^V$ is
 at most
 \[
  \frac{2^{t_V}R_U}{\prod_i(1-e_{t_i}(p,n))}\ \le\ \frac{2^{t_V}R_U}{(1-e_\Delta(p,n))^{|V|}} .
 \]
\end{enumerate}
\end{lemma}

\begin{proof}
\ref{it:vertex:inv} is \cref{lem:marginals}, parts (1) and (2): the map
$z\mapsto\alpha z+\beta$ is a bijection of $\operatorname{Inj}_p(G)$.

\ref{it:vertex:dens} For fixed $z_{W\setminus\{w\}}=y'$ the elements of
$\operatorname{Inj}_p(G)$ extending $y'$ correspond to the admissible values of
$z_w$, all equally likely under $\lambda$. They form a set $\Phi$ of
\cref{lem:extension} whose roots are the $d_w$ (distinct) labels of the
neighbours of $w$ and whose excluded set consists of the other $|W|-1\le n$
labels, so \cref{lem:extension}\ref{it:ext:coarse} gives the lower bound. Under
$m$ the conditional probability of $a\in\Phi_w$ is
$(dm/d\lambda)(y',a)/\sum_{a'\in\Phi_w}(dm/d\lambda)(y',a')\le Rc/(|\Phi_w|c)$.

\ref{it:vertex:subset} The marginal of $\lambda$ at $x\in\Fp^J$ is
$C(x)/|\operatorname{Inj}_p(G)|$, $C(x)$ the number of completions. By
\cref{lem:graphs:sequential}\ref{it:seq:count},\ref{it:seq:coarse} for the anchors
$J$ and any ordering of $W\setminus J$, whose numbers $t_i$ of earlier neighbours
satisfy $t_i\le d_i$, $C$ is positive exactly on $\operatorname{Inj}_p(G[J])$ and
takes values in an interval whose endpoints have ratio at most
$\prod(1+e_{t_i})/(1-e_{t_i})\le R_{W\setminus J}$ (by \ref{it:et:poly}). The density with
respect to the uniform law on $\operatorname{Inj}_p(G[J])$ is proportional to $C$.
For $m$, $m_J(x)$ is the sum of $m$ over the completions of $x$, which lies
between $cC(x)/|\operatorname{Inj}_p(G)|$ and $RcC(x)/|\operatorname{Inj}_p(G)|$.

\ref{it:vertex:del} is \ref{it:vertex:subset} with $J=W\setminus\{w\}$.

\ref{it:vertex:tuple} The conditional probability of $x_V$ is
$C(x_C,x_V)/\sum_{x'_V}C(x_C,x'_V)$, where $C(x_C,x_V)$ is the number of completions
of $(x_C,x_V)$ on $U$. By \cref{lem:graphs:sequential}\ref{it:seq:count} with the
anchors $C\cup V$, this count is positive exactly when
$(x_C,x_V)\in\operatorname{Inj}_p(G[C\cup V])$, and it takes values in an interval
with endpoint ratio at most $R_U$; so the admissible tuples are injective, and
the conditional probability of each is at most $R_U$ divided by their number.
Since $x_C$ is feasible, $x_C\in\operatorname{Inj}_p(G[C])$, and
\cref{lem:graphs:sequential}\ref{it:seq:count} for $G[C\cup V]$ with the anchors $x_C$
and the given ordering of $V$ shows that there are at least
$p^{|V|}2^{-t_V}\prod_i(1-e_{t_i}(p,n))$ admissible tuples. Multiplying by
$p^{|V|}$ gives the bound, and the last inequality is \ref{it:et:poly} with
$t_i\le\Delta$.
\end{proof}

For a block $\Blk$ all these statements hold at each prime of $\Blk$, and the
corresponding statements for $\lambda_{\Blk}(G)$ follow because its marginals,
conditional laws and densities are products over the primes of $\Blk$.

\subsection{Comparing two graph laws: a ratio of four counts}\label{subsec:labellings:four-count}

In Part~V (\cref{lem:front12:reference,lem:upper-reference-rp}) the conditional laws given a middle label (the rows of the middle and upper references) of two graph laws whose graphs differ only away from the conditioned and queried
vertices are compared. Their density is then a ratio of four completion counts, whose main
terms cancel.

\begin{lemma}[Four-count comparison]\label{lem:four-count}
Let $G$ and $G'$ be oriented graphs with vertex sets $W$ and $W'$, and let
$M,S$ be disjoint subsets of $W\cap W'$ such that $G[M\cup S]=G'[M\cup S]$ (the
same edges with the same orientations). Fix orderings and one-step error pairs
for which \cref{lem:graphs:sequential} applies to $G$ with the anchors $M$ and with
the anchors $M\cup S$, and to $G'$ with the same anchors, and write
$\mathrm{up}_G(A),\mathrm{lo}_G(A)$, $\mathrm{up}_{G'}(A),\mathrm{lo}_{G'}(A)$ for
the logarithmic count errors with the anchors $A$; write $C_G^A(x_A)$ for the
number of elements of $\operatorname{Inj}_p(G)$ extending $x_A$, so that
$C_G^\emptyset=|\operatorname{Inj}_p(G)|$. Let $\lambda=\lambda_p(G)$ and $\lambda'=\lambda_p(G')$.
\begin{enumerate}[label=\textup{(\alph*)}]
\item\label{it:four:formula} For every $x_M\in\operatorname{Inj}_p(G[M])$, the conditional laws
 of $z_S$ given $z_M=x_M$ under $\lambda$ and under $\lambda'$ are both supported on
 $\Sigma(x_M)=\{x_S:(x_M,x_S)\in\operatorname{Inj}_p(G[M\cup S])\}$, and the density
 of the first with respect to the second is
 \[
  \varrho(x_M,x_S)=\frac{C_G^{M\cup S}(x_M,x_S)\,C_{G'}^M(x_M)}
  {C_G^M(x_M)\,C_{G'}^{M\cup S}(x_M,x_S)} .
 \]
\item\label{it:four:errors} The main terms cancel, and
 \begin{align*}
  \log\varrho&\le\mathrm{up}_G(M\cup S)+\mathrm{lo}_G(M)+\mathrm{lo}_{G'}(M\cup S)
  +\mathrm{up}_{G'}(M),\\
  -\log\varrho&\le\mathrm{lo}_G(M\cup S)+\mathrm{up}_G(M)+\mathrm{up}_{G'}(M\cup S)
  +\mathrm{lo}_{G'}(M).
 \end{align*}
 With $M=\emptyset$ this compares the marginals of $\lambda$ and $\lambda'$ on $S$.
\item\label{it:four:block} For a block $\Blk$, the density of the conditional law of
 $z_S$ given $z_M$ under $\lambda_{\Blk}(G)$ with respect to that under
 $\lambda_{\Blk}(G')$ is the product over $p\in\Blk$ of the densities of
 \ref{it:four:formula}, and its logarithm lies in the sum over $p\in\Blk$ of the
 intervals of \ref{it:four:errors}.
\end{enumerate}
\end{lemma}

\begin{proof}
\ref{it:four:formula} Under $\lambda$ the law of $z_{M\cup S}$ is
$C_G^{M\cup S}/|\operatorname{Inj}_p(G)|$ and that of $z_M$ is
$C_G^M/|\operatorname{Inj}_p(G)|$, so the conditional law of $z_S$ given $z_M=x_M$
is $C_G^{M\cup S}(x_M,\cdot)/C_G^M(x_M)$. By \cref{lem:graphs:sequential}\ref{it:seq:count}
it is positive exactly on $\Sigma(x_M)$. The same holds for $G'$, with the same
set $\Sigma(x_M)$ because $G'[M\cup S]=G[M\cup S]$, and the ratio is $\varrho$.

\ref{it:four:errors} For a graph $H$ with vertex set $W_H$ write
$\mathrm{main}_H(A)=p^{|W_H|-|A|}2^{-(|E(H)|-|E(H[A])|)}$ for the main term of
\cref{lem:graphs:sequential}. In the quotient of the four main terms of $\varrho$
the powers of $p$ cancel, and the powers of $2$ combine to
$2^{|E(G[M\cup S])|-|E(G[M])|-|E(G'[M\cup S])|+|E(G'[M])|}=1$, since $G$ and $G'$
have the same induced graphs on $M\cup S$ and on $M$. By
\cref{lem:graphs:sequential}\ref{it:seq:count} each count $C_H^A$ lies in
$\mathrm{main}_H(A)\,[e^{-\mathrm{lo}_H(A)},e^{\mathrm{up}_H(A)}]$; inserting
these four intervals gives the bounds.

\ref{it:four:block} The labels at different primes are independent under both
product laws, so the conditional laws and their density are products.
\end{proof}

In Part~V the lemma is applied to the field graphs $G_F$ and $G_E$ of a fold
record and of a matching record on three objects (the fold and matching patterns of \cref{def:patterns}), whose induced graphs on any
two adjacent objects coincide (\cref{lem:layer12-geometry}(d) in the middle
range, \cref{lem:upper-geometry}(d) in the upper range): on two adjacent
objects the fold relation and the matching relation both relate every pair of
distinct object values and neither has a forward pair. With $M$ the middle
object and $S$ one of the outer objects, part \ref{it:four:formula} writes the
density of the conditional outer marginals of the rows as a ratio of
completion counts, and part \ref{it:four:errors} bounds its logarithm: in the
middle range the counts are exact (\cref{lem:record:class-three}(iii)), and in
the upper range part \ref{it:four:errors} with the graded insertion counts gives
the pointwise marginal ratio of the rows \ref{it:record:upper-F5}. With
$M=\emptyset$ and $S$ the middle object, part \ref{it:four:errors} bounds the
density of the middle marginals at first order. On a block whose primes are at
least $b$ and whose logarithmic length $\sum\log p$ is at most
$L$ (any such bound; $L=570$ in the construction), the errors,
which are
nondecreasing in $r$ by \cref{lem:graphs:sequential}\ref{it:seq:thresholds}, add
over at most $L/\log b$ primes by part \ref{it:four:block}.

\subsection{How uniform is one label, or a queried pair, given the others?}\label{subsec:labellings:trace}

The mixing hypothesis (R4) of \cref{hyp:record:reference} needs, at large primes, that conditioning on all other
slots leaves the maximal correlation of one slot, or of a queried pair of slots, with the rest small (\cref{def:mixing}). For a graph law this follows from the symmetry of
$\lambda_p(G)$ alone: the operator $T^*T$ commutes with translations and with
square dilations, so it is diagonal, or block diagonal, in the additive
characters with equal traces along the dilation orbits, and its trace is an
average of reciprocal counts.

\begin{lemma}[Trace lemma for conditional mixing]\label{lem:trace-mixing}
Let $G=(W,E)$ be an oriented graph with $n=|W|\ge3$ vertices and maximum degree
$\Delta$, let $p$ be an odd prime with $\operatorname{Inj}_p(G)\ne\emptyset$, and
let $\lambda=\lambda_p(G)$.
\begin{enumerate}[label=\textup{(\alph*)}]
\item\label{it:trace:slot} \emph{One slot.} Let $v\in W$, $d=\deg_G(v)$, and let
 $w<1$ with $w\ge e_d(p,n-1)$ (for instance $w=w_\Delta(p,n)$ if this is below
 one). The conditional
 operator $T=T^{\lambda,v}$ of \cref{def:mixing} satisfies, for every
 $f\colon\Fp\to\C$ with $\E f=0$,
 \[
  \|Tf\|^2\le\frac2{p-1}\Bigl(\frac{2^d}{1-w}-1\Bigr)\|f\|^2 .
 \]
\item\label{it:trace:pair} \emph{A pair, trace form.} Let $(u,v)\in E$ and
 $W'=W\setminus\{u,v\}$. The pair $(z_u,z_v)$ has law $\Unif(\Om)$. For a
 labelling $z'$ of $W'$ in the support of the marginal of $\lambda$, let
 $\Phi(z')\subseteq\Om$ be the set of pairs $(x,x')$ such that $z'$ together with
 $z_u=x$, $z_v=x'$ lies in $\operatorname{Inj}_p(G)$; for $y\in\Qsq$ let
 $c_y(z')$ be the number of $x$ with $(x,x+y)\in\Phi(z')$, and put
 $\varphi_y(z')=\frac{p-1}2\,c_y(z')/|\Phi(z')|-1$. The conditional operator
 $(T_{uv}f)(z')=\E_\lambda[f(z_u,z_v)\mid z_{W'}=z']$ satisfies, for every
 $f\in L^2(\Unif(\Om))$ with mean zero,
 \[
  \|T_{uv}f\|^2\le\max\Bigl\{\E_\lambda\Bigl[\frac2{p-1}\sum_{y\in\Qsq}\varphi_y^2\Bigr],\
  p\,\E_\lambda\Bigl[\frac1{|\Phi|}\Bigr]\Bigr\}\|f\|^2 .
 \]
\item\label{it:trace:explicit} \emph{A pair, explicit form.} In \ref{it:trace:pair}
 let $d'_u=\deg_G(u)-1$, $d'_v=\deg_G(v)-1$, $t=d'_u+d'_v$, $s=\max(d'_u,d'_v)$,
 $n_\times=d'_ud'_v$, and let $e<1$ with $e\ge e_t(p,2n-4)$ and $e\ge e_{s+1}(p,n)$.
 Put
 \[
  \eta=\frac{n_\times2^{t-s+1}}{p-1},\qquad
  \varphi_{\rm gen}=\frac{3e}{(1-e)^2}+(1+e)\eta,\qquad
  \varphi_{\rm exc}=\frac{2^{t-s}(1+e)}{(1-e)^2} .
 \]
 Then for every mean-zero $f\in L^2(\Unif(\Om))$,
 \[
  \|T_{uv}f\|^2\le\max\Bigl\{\varphi_{\rm gen}^2+\frac{2n_\times}{p-1}\varphi_{\rm exc}^2,\
  \frac{2^{t+1}}{(p-1)(1-e)^2}\Bigr\}\|f\|^2 .
 \]
\item\label{it:trace:72} \emph{Seventy-two vertices of degree at most twenty.} If
 $n=72$ and $\Delta\le20$, then for every prime $p\ge2^{512}$, every vertex $v$ and
 every edge $(u,v)$, the squared norms in \ref{it:trace:slot} and
 \ref{it:trace:pair} on the mean-zero functions are at most $2^{22}/p$ and
 $2^{92}/p$ respectively.
\end{enumerate}
\end{lemma}

\begin{proof}
Let $\mathfrak G$ be the group of the maps $x\mapsto\alpha x+\beta$ ($\alpha\in\Qsq$,
$\beta\in\Fp$), acting on labellings coordinatewise. It preserves
$\operatorname{Inj}_p(G)$ (it multiplies every difference by a nonzero square),
hence $\lambda$ and all its marginals (\cref{lem:marginals}); its translations act
transitively on $\Fp$, and $\mathfrak G$ acts transitively on $\Om$
(\cref{lem:marginals}, part (3)), so the one-slot marginals are uniform and, the
edge $(u,v)$ forcing $(z_u,z_v)\in\Om$, the pair has law $\Unif(\Om)$.

\ref{it:trace:slot} For a labelling $z'$ of $W\setminus\{v\}$ in the support of the
marginal $\lambda'$, let $\Phi_v(z')$ be the set of admissible labels of $v$; since
$\lambda$ is uniform, $Tf(z')=|\Phi_v(z')|^{-1}\sum_{x\in\Phi_v(z')}f(x)$. The labels
of the $d$ neighbours of $v$ are distinct and lie among the $n-1$ excluded
labels, so \cref{lem:extension}\ref{it:ext:coarse} gives
$|\Phi_v(z')|\ge p2^{-d}(1-e_d(p,n-1))\ge p2^{-d}(1-w)$. For $g\in\mathfrak G$,
$\Phi_v(gz')=g\Phi_v(z')$, hence $T(f\circ g)=(Tf)\circ g$; the maps $f\mapsto f\circ g$
are unitary on $L^2(\Fp)$ and on $L^2(\lambda')$, so $S=T^*T$ commutes with them.
Commuting with the translations, $S$ is diagonal in the characters
$e_\xi(x)=e(\xi x/p)$, each of which spans the joint eigenspace of the
translations for its own character: $Se_\xi=\kappa_\xi e_\xi$ with
$\kappa_\xi=\|Te_\xi\|^2\ge0$. Since $e_\xi(\alpha x)=e_{\alpha\xi}(x)$, commuting with
the dilations gives $\kappa_{\alpha\xi}=\kappa_\xi$ for $\alpha\in\Qsq$: $\kappa$ takes one
value $\kappa_+$ on $\Qsq$ and one value $\kappa_-$ on the nonsquares, and
$\kappa_0=\|T\one\|^2=1$. In the orthonormal basis $(\sqrt p\,\one_{\{x\}})_x$,
\[
 \tr S=\sum_xp\,\|T\one_{\{x\}}\|^2
 =\sum_xp\,\E_{\lambda'}\Bigl[\frac{\one\{x\in\Phi_v\}}{|\Phi_v|^2}\Bigr]
 =p\,\E_{\lambda'}\Bigl[\frac1{|\Phi_v|}\Bigr]\le\frac{2^d}{1-w},
\]
and $\tr S=1+\frac{p-1}2(\kappa_++\kappa_-)$. Hence
$\max(\kappa_+,\kappa_-)\le\frac2{p-1}\bigl(\frac{2^d}{1-w}-1\bigr)$, and for
$f=\sum_{\xi\ne0}\widehat f(\xi)e_\xi$,
$\|Tf\|^2=\sum_{\xi\ne0}\kappa_\xi|\widehat f(\xi)|^2\le\max(\kappa_+,\kappa_-)\|f\|^2$.
Finally $e_d(p,n-1)\le e_\Delta(p,n)\le w_\Delta(p,n)$ by \ref{it:et:compare},
\ref{it:et:poly}.

\ref{it:trace:pair} Use the coordinates $(x,y)$, $x'=x+y$, which identify $\Om$
with $\Fp\times\Qsq$ and $\Unif(\Om)$ with the uniform law there; write
$\E_y$ for the average over $y\in\Qsq$ and $\lambda''$ for the marginal of $z_{W'}$.
Since $\lambda$ is uniform, $T_{uv}f(z')=|\Phi(z')|^{-1}\sum_{\omega\in\Phi(z')}f(\omega)$, and
as in \ref{it:trace:slot}, $\Phi(gz')=g\Phi(z')$ for $g\in\mathfrak G$ (acting diagonally
on $\Om$), so $S=T_{uv}^*T_{uv}$ commutes with the unitaries $f\mapsto f\circ g$. The
translation $x\mapsto x+\beta$ acts by $(x,y)\mapsto(x+\beta,y)$ and the dilation by
$\alpha\in\Qsq$ by $(x,y)\mapsto(\alpha x,\alpha y)$. For $k\in\Fp$ let
$\mathcal H_k=\{e_k(x)\psi(y)\}$; these spaces are orthogonal, they are the joint
eigenspaces of the translations (with distinct characters), and
$L^2_0(\Unif(\Om))=(\mathcal H_0\ominus\C\one)\oplus\bigoplus_{k\ne0}\mathcal H_k$. So $S$
preserves each $\mathcal H_k$, and also $\mathcal H_0\ominus\C\one$ (it fixes $\one$
and is self-adjoint), and $\|T_{uv}|_{L^2_0}\|^2$ is the largest of the norms of
the restrictions $S_k$ of $S$ to $\mathcal H_k$ ($k\ne0$) and to
$\mathcal H_0\ominus\C\one$.

\emph{Nonzero frequencies.} The dilation by $\alpha$ maps $\mathcal H_k$ onto
$\mathcal H_{\alpha k}$ and commutes with $S$, so $\tr S_k$ takes one value
$\sigma_+$ for $k\in\Qsq$ and one value $\sigma_-$ for the nonsquares. Since
$S_k\succeq0$, $\|S_k\|\le\tr S_k$, and
$\frac{p-1}2(\sigma_++\sigma_-)\le\tr S$. In the orthonormal basis
$(|\Om|^{1/2}\one_{\{\omega\}})_\omega$, as in \ref{it:trace:slot},
$\tr S=|\Om|\,\E_{\lambda''}[1/|\Phi|]$. With $|\Om|=p(p-1)/2$ this gives
$\|S_k\|\le\sigma_++\sigma_-\le p\,\E[1/|\Phi|]$.

\emph{Frequency zero.} For $\psi\in\mathcal H_0$ with $\E_y\psi=0$,
\[
 T_{uv}\psi(z')=\sum_{y\in\Qsq}\frac{c_y(z')}{|\Phi(z')|}\psi(y)
 =\E_y\bigl[(1+\varphi_y(z'))\psi(y)\bigr]=\E_y\bigl[\varphi_y(z')\psi(y)\bigr],
\]
so by Cauchy--Schwarz
$|T_{uv}\psi(z')|^2\le\E_y[\varphi_y(z')^2]\,\E_y|\psi|^2$, and
$\|T_{uv}\psi\|^2\le\E_{\lambda''}\bigl[\frac2{p-1}\sum_y\varphi_y^2\bigr]\|\psi\|^2$.

\ref{it:trace:explicit} Fix $z'$ in the support; let $\mathcal L=z'(W')$, a set of $n-2$
labels, and let $B_u$, $B_v$ be the sets of labels of the neighbours of $u$,
respectively $v$, in $W'$; $|B_u|=d'_u$, $|B_v|=d'_v$. A pair $(x,x+y)$ lies in
$\Phi(z')$ iff $x,x+y\notin\mathcal L$ and $x$ satisfies the $d'_u$ conditions
$\chi(x-\beta)=s_\beta$ ($\beta\in B_u$) from the edges at $u$ and the $d'_v$
conditions $\chi(x-(\beta'-y))=s'_{\beta'}$ ($\beta'\in B_v$) from the edges at $v$; the
edge $(u,v)$ is the condition $y\in\Qsq$. So $c_y$ is the size of a set $\Phi$ of
\cref{lem:extension} with excluded set $\mathcal L\cup(\mathcal L-y)$, of size at most $2n-4$,
which contains all the roots. Call $y$ \emph{generic} if
$B_u\cap(B_v-y)=\emptyset$; at most $n_\times$ values of $y$ are not.

\emph{Counts.} For generic $y$ there are $t$ distinct roots, and
\cref{lem:extension}\ref{it:ext:coarse} gives $c_y\in p2^{-t}[1-e,1+e]$. For every
$y$, keeping only the conditions of an endpoint with $s$ conditions,
\cref{lem:extension}\ref{it:ext:upper} gives $c_y\le p2^{-s}(1+e^+_s)\le p2^{-s}(1+e)$,
since $e^+_s\le e_s(p,n)\le e_{s+1}(p,n)\le e$. Labelling first the endpoint
with $s$ conditions and then the other endpoint, which has $t-s+1\le s+1$
conditions (its own and the edge $(u,v)$),
\cref{lem:graphs:sequential}\ref{it:seq:count} with the coarse pair gives
$|\Phi(z')|\ge p^22^{-(t+1)}(1-e)^2$.

\emph{Bounds.} Hence $p\,\E[1/|\Phi|]\le2^{t+1}/(p(1-e)^2)$, which is at most the
second term. Recall $|\Qsq|=\frac{p-1}2$. For generic $y$,
$|\Qsq|c_y/|\Phi|\le|\Qsq|p2^{-t}(1+e)/(p^22^{-(t+1)}(1-e)^2)\le(1+e)/(1-e)^2$, so
$\varphi_y\le3e/(1-e)^2$; and
$|\Phi|=\sum_yc_y\le|\Qsq|p2^{-t}(1+e)+n_\times p2^{-s}(1+e)=|\Qsq|p2^{-t}(1+e)(1+\eta)$,
so $|\Qsq|c_y/|\Phi|\ge(1-e)/((1+e)(1+\eta))$ and
\[
 -\varphi_y\le\frac{2e+(1+e)\eta}{(1+e)(1+\eta)}\le2e+(1+e)\eta .
\]
Thus $|\varphi_y|\le\varphi_{\rm gen}$ for generic $y$. For the other $y$,
$-1\le\varphi_y\le|\Qsq|p2^{-s}(1+e)/(p^22^{-(t+1)}(1-e)^2)\le\varphi_{\rm exc}$, and
$\varphi_{\rm exc}\ge1$, so $|\varphi_y|\le\varphi_{\rm exc}$. Therefore
$\frac2{p-1}\sum_y\varphi_y^2\le\varphi_{\rm gen}^2+\frac{2n_\times}{p-1}\varphi_{\rm exc}^2$,
and \ref{it:trace:explicit} follows from \ref{it:trace:pair}.

\ref{it:trace:72} In \ref{it:trace:slot} take $w=w_{20}(p,72)=2^{20}(10r+72r^2)$,
which is at least $e_d(p,71)$ for $d\le20$ by \ref{it:et:compare},
\ref{it:et:poly}; the bound is nondecreasing in $d$, so $d=20$ is the worst case.
In \ref{it:trace:explicit} take $e=e_{38}(p,140)=2^{38}((18+2^{-38})r+140r^2)$,
which is at least $e_t(p,140)$ and $e_{s+1}(p,72)$ since $t\le38$ and $s+1\le20$;
the right side of \ref{it:trace:explicit} is nondecreasing in $n_\times\le361$,
$t-s\le19$ and $t+1\le39$, so the worst case is $d'_u=d'_v=19$, with
$\eta=361\cdot2^{20}/(p-1)$ and $\varphi_{\rm exc}=2^{19}(1+e)/(1-e)^2$. Multiplied by
$p$, the three resulting expressions are
\[
 \frac2{1-r^2}\Bigl(\frac{2^{20}}{1-w}-1\Bigr),\qquad
 \Bigl(\frac{\varphi_{\rm gen}}r\Bigr)^2+\frac{722\,\varphi_{\rm exc}^2}{1-r^2},\qquad
 \frac{2^{39}}{(1-r^2)(1-e)^2},
\]
and $e/r$, $w/r$ and $\eta/r=361\cdot2^{20}r/(1-r^2)$ are nondecreasing in $r$; so
all three are nondecreasing in $r$ (\cref{lem:majorant}), and it suffices to
bound them at $p=2^{512}$, which \cref{cert:trace-mixing-72} does.
\end{proof}

\begin{certificate}[Mixing constants of the two-layer root at $p=2^{512}$]
\label{cert:trace-mixing-72}
At $p=2^{512}$, the quantities $e=e_{38}(p,140)$,
$\eta=361\cdot2^{20}/(p-1)$, $\varphi_{\rm gen}=3e/(1-e)^2+(1+e)\eta$,
$\varphi_{\rm exc}=2^{19}(1+e)/(1-e)^2$ and $w=2^{20}(10p^{-1/2}+72p^{-1})$
of the proof of \cref{lem:trace-mixing}\ref{it:trace:72} satisfy
$e<\frac12$, $w<\frac12$ and the three inequalities below. They are
evaluated in exact rational arithmetic by the function
\texttt{leading\_trace} of
\path{research/push_2026_10_17/structural/layer72_analytic.py}:
\[
 p\Bigl(\varphi_{\rm gen}^2+\frac{722}{p-1}\varphi_{\rm exc}^2\Bigr)<2^{92},\qquad
 \frac{p\,2^{39}}{(p-1)(1-e)^2}<2^{92},\qquad
 \frac{2p}{p-1}\Bigl(\frac{2^{20}}{1-w}-1\Bigr)<2^{22} .
\]
For orientation only (not part of the check), the three left sides are approximately $2^{87.51}$, $2^{39}$ and $2^{21}$. How the check
is run and authenticated is recorded in its row of \cref{tab:cert:finite}.
\end{certificate}

\subsection{Primes $p\equiv3\pmod4$: the Paley tournament, triples, and operators}\label{subsec:labellings:class-three}

At a prime $p\equiv3\pmod4$, $-1$ is a nonsquare, so for $x\ne y$ exactly one of
$y-x$ and $x-y$ is a square: the strict-square relation is a tournament, the
Paley tournament, and the character kernel $\chi(t-s)$ is antisymmetric. This is
why the middle and upper primes are taken with
$p\equiv3\pmod4$: the laws of directed
configurations there have exact counts and exact correlation identities, which
the following lemma collects. (At $p\equiv1\pmod4$ the relation is symmetric
and the tournament statements of the lemma fail; its part \ref{it:paley:corr}
holds for every odd $p$, and the disjoint-triple operator at the end of the
subsection uses no character sum and is treated at every prime $p\ge7$.)

\begin{lemma}[Paley-triple calculus]\label{lem:paley-triple}
Let $p\equiv3\pmod4$ be prime and write $a\to b$ if $\chi(b-a)=1$. Put
$k=(p-3)/4$ and let $P$ be the Paley operator of \cref{lem:square-shift},
$(Ph)(s)=\E_t\chi(t-s)h(t)$ on $L^2(\Fp)$.
\begin{enumerate}[label=\textup{(\alph*)}]
\item\label{it:paley:counts} \emph{Tournament counts.} $\to$ is a tournament on $\Fp$
 in which every label has $(p-1)/2$ out-neighbours and $(p-1)/2$
 in-neighbours. Two distinct labels $a\ne a'$ have exactly $k$ common
 out-neighbours and exactly $k$ common in-neighbours, and
 $\#\{c:a\to c\to a'\}=k+1-\one[a\to a']$.
\item\label{it:paley:corr} \emph{Character correlations.}
 $\E_t\chi(t-a)\chi(t-a')=\one[a=a']-q$ for all $a,a'\in\Fp$ (for every odd $p$;
 this is \eqref{eq:two-root}).
\item\label{it:paley:op} \emph{The Paley operator.} $P^*=-P$, $P\one=0$,
 $P^*P=PP^*=q(I-\Jproj)$ and $\|P\|=r$ (\cref{lem:square-shift}(c)). Hence $P$
 maps $H_0$ into itself, and on $H_0$ we have $P^2=-qI$, i.e.\ $pP^2=-I$, and
 $P/r$ is unitary with eigenvalues $\pm i$.
\item\label{it:paley:triples} \emph{Directed triples.} Let $\nu_F$ be the uniform law on
 $H_F=\{(x_0,x_1,x_2):\ x_0,x_1,x_2\ \text{distinct},\ x_0\to x_2\}$, and
 $V_i\colon L^2(\Fp)\to L^2(\nu_F)$, $(V_ih)(x)=h(x_i)$. The coordinates of
 $\nu_F$ are uniform, each $V_i$ is an isometry, and $V_i^*$ is the conditional
 expectation onto $x_i$. On $H_0$,
 \[
  V_0^*V_2=\frac{P-qI}{1-q},\qquad
  V_0^*V_1=V_1^*V_2=-\frac{I+V_0^*V_2}{p-2}=-\frac I{p-1}-\frac{pP}{(p-1)(p-2)},
 \]
 and so $V_2^*V_0=-(P+qI)/(1-q)$ and
 $V_1^*V_0=V_2^*V_1=-\frac I{p-1}+\frac{pP}{(p-1)(p-2)}$. The norms on $H_0$ are
 exactly
 \[
  \|V_0^*V_2\|=\varrho:=\frac{r\sqrt{1+q}}{1-q}=\frac{\sqrt{p+1}}{p-1},\qquad
  \|V_0^*V_1\|=\|V_1^*V_2\|=\varrho_1:=\frac{q}{1-q}\sqrt{1+\frac q{(1-2q)^2}},
 \]
 and $\varrho_1\le\varrho':=q(1+\varrho)/(1-2q)$. For $p\ge7$, $\varrho'\le\varrho$, and
 the Gram operator $(V_i^*V_j)_{i,j}$ of $(V_0,V_1,V_2)$ on $H_0^3$ has norm at
 most $\Gamma:=1+\varrho+\varrho'$.
\item\label{it:paley:features} \emph{Features.} Let $A_1=V_2P$, $A_2=V_1P$ and
 $(A_3h)(x)=-\E_s[\chi(s-x_2)\chi(s-x_1)h(s)]-q\,\E h$, operators
 $L^2(\Fp)\to L^2(\nu_F)$. Each $A_j$ annihilates the constants and maps into the
 $\nu_F$-mean-zero functions, and on $H_0$
 \begin{gather*}
  A_1^*A_1=A_2^*A_2=qI,\qquad A_3^*A_3=\frac{1-2q-q^2}{p-1}I\preceq qI,\qquad
  A_1^*A_2=-\frac q{p-1}I+\frac P{(p-1)(p-2)},\\
  A_1^*A_3=-\frac q{(p-1)(p-2)}I+\frac q{p-1}P,\qquad
  A_2^*A_3=\frac q{(p-1)(p-2)}I+\frac q{p-1}P,\\
  V_1^*A_3=-\frac q{p-1}I+\frac P{(p-1)(p-2)},\qquad
  V_2^*A_3=-\frac q{p-1}I-\frac P{(p-1)(p-2)},\\
  V_0^*A_1=-\frac{I+P}{p-1},\qquad V_0^*A_2=\frac I{(p-1)(p-2)}-\frac P{p-1},\qquad
  V_0^*A_3=\frac{2I}{p(p-1)(p-2)} .
 \end{gather*}
 Hence, for $\xi\in\R^3$ and $A_\xi=\xi_1A_1+\xi_2A_2+\xi_3A_3$, $V_0^*A_\xi=a_\xi I+b_\xi P$
 on $H_0$ with
 \[
  a_\xi=-\frac{\xi_1q}{1-q}+\frac{\xi_2q^2}{(1-q)(1-2q)}+\frac{2\xi_3q^3}{(1-q)(1-2q)},
  \qquad b_\xi=-\frac{(\xi_1+\xi_2)q}{1-q},
 \]
 and $\|(A_1,A_2,A_3)\|\le r\gamma$ on $H_0^3$, with $\gamma^2=1+q(1+2r)/(1-2q)$.
\item\label{it:paley:perturb} \emph{Density perturbation.} Let $m=g\nu_F$ be a
 probability measure on $H_F$ with uniform coordinate marginals and
 $|g-1|\le e$. Then $|\ip FG_m-\ip FG_{\nu_F}|\le e\|F\|_{\nu_F}\|G\|_{\nu_F}$ and
 $(1-e)\|F\|^2_{\nu_F}\le\|F\|_m^2\le(1+e)\|F\|^2_{\nu_F}$ for all $F,G$. In
 particular, for $f,h\in H_0$ the correlations $\ip{V_if}{V_jh}$ under $m$ and
 under $\nu_F$ differ by at most $e\|f\|\|h\|$, the norms $\|V_if\|$ being equal
 under both laws; and for every operator $A\colon L^2(\Fp)\to L^2(\nu_F)$, the
 adjoints of $V_i$ in $L^2(m)$ and in $L^2(\nu_F)$ satisfy
 $\|V_{i,m}^*A-V_{i,\nu_F}^*A\|\le e\|A\|$.
\end{enumerate}
\end{lemma}

\begin{proof}
\ref{it:paley:corr} is \eqref{eq:two-root}.

\ref{it:paley:counts} Since $\chi(-1)=-1$, $\chi(a-b)=-\chi(b-a)$, so exactly one of
$a\to b$, $b\to a$ holds for $a\ne b$; out- and in-neighbours of $a$ are the
$(p-1)/2$ labels $b$ with $\chi(b-a)=1$, respectively $-1$. For $a\ne a'$,
\[
 4\#\{c:a\to c,\ a'\to c\}=\sum_{c\ne a,a'}(1+\chi(c-a))(1+\chi(c-a'))
 =(p-2)-\chi(a'-a)-\chi(a-a')-1=p-3,
\]
using $\sum_{c\ne a}\chi(c-a)=0$ and \eqref{eq:two-root}; the in-neighbours are the
same computation with $1-\chi$, and give $p-3$ again. For the paths,
$4\#\{c\}=\sum_{c\ne a,a'}(1+\chi(c-a))(1+\chi(a'-c))=(p-2)-2\chi(a'-a)+1$,
because both linear sums equal $-\chi(a'-a)$ and
$\sum_c\chi(c-a)\chi(a'-c)=-\sum_c\chi((c-a)(c-a'))=1$; this is $4k$ if $a\to a'$
and $4k+4$ otherwise.

\ref{it:paley:op} The first assertions are \cref{lem:square-shift}(c), since
$\chi(-1)=-1$. As $P^*\one=-P\one=0$, $P$ maps $H_0$ into itself; there
$P^*P=PP^*=qI$, so $P/r$ is unitary, and $P^2=-P^*P=-qI$, so its eigenvalues
are $\pm i$.

\ref{it:paley:triples} The translations preserve $H_F$ and act transitively on
each coordinate, so the coordinates are uniform, $V_i$ is an isometry, and
$\ip{V_ih}F_{\nu_F}=\E_s\overline{h(s)}\,\E_{\nu_F}[F\mid x_i=s]$. Given $x_0=s$, the label
$x_2$ is uniform on the set $s+\Qsq$ of the out-neighbours of $s$, and given
$(x_0,x_2)$ the label $x_1$ is uniform on the other $p-2$ labels. So $V_0^*V_2$ is
the square shift $T_p$ of \cref{lem:square-shift}, which is
$(p\Jproj-I+pP)/(p-1)$ by part (c) there, that is $(pP-I)/(p-1)$ on $H_0$; and
for $g\in H_0$, $\E[g(x_1)\mid x_0,x_2]=-(g(x_0)+g(x_2))/(p-2)$, which gives $V_0^*V_1$;
conditioning on $x_2$ instead gives $V_2^*V_1=-(I+V_2^*V_0)/(p-2)$, whose adjoint
is $V_1^*V_2$. The closed forms follow from $q/(1-q)=1/(p-1)$, and the adjoints
from $P^*=-P$. By \ref{it:paley:op}, $P$ is normal on $H_0$ with eigenvalues
$\pm ir$, so $V_0^*V_2$ has eigenvalues $(-q\pm ir)/(1-q)$, of modulus $\varrho$, and
$V_0^*V_1$ has eigenvalues $-\frac1{p-1}\mp\frac{ipr}{(p-1)(p-2)}$, of modulus
$\frac1{p-1}(1+\frac p{(p-2)^2})^{1/2}=\varrho_1$; and
$\varrho_1=\|I+V_0^*V_2\|/(p-2)\le(1+\varrho)/(p-2)=\varrho'$. The inequality $\varrho'\le\varrho$
is $q\le\varrho(1-3q)$, i.e.\ $1\le(p-3)\sqrt{p+1}/(p-1)$, which holds for $p\ge7$.
The Schur test applied to the symmetric matrix of the norms of the blocks of the
Gram operator bounds its norm by the largest row sum,
$1+\max(\varrho+\varrho_1,2\varrho_1)\le\Gamma$.

\ref{it:paley:features} By \ref{it:paley:op}, $A_1\one=A_2\one=0$, and
$A_3\one=-(\one[x_1=x_2]-q)-q=0$ by \ref{it:paley:corr}. The functions $A_1h$, $A_2h$
have mean $\E Ph=0$, and $\E_{\nu_F}[\chi(s-x_2)\chi(s-x_1)]$ does not depend on $s$
(translation invariance), hence equals its average over $s$, which is $-q$; so
$\E_{\nu_F}A_3h=q\E h-q\E h=0$.

Let $M_{1|2}=V_2^*V_1$, the conditional expectation of a function of $x_1$ given
$x_2$. For distinct $a,b$ the number of $x_0$ with $(x_2,x_1)=(a,b)$ is
$(p-2-\chi(a-b))/2$; summing $g(b)$ against these weights gives, on all of
$L^2(\Fp)$,
\[
 M_{1|2}=\frac{p\Jproj-I}{p-1}+\frac{pP}{(p-1)(p-2)},\qquad
 M_{1|2}^*=V_1^*V_2=\frac{p\Jproj-I}{p-1}-\frac{pP}{(p-1)(p-2)} .
\]
On $H_0$, where $\Jproj=0$ and $P$ maps $H_0$ into itself,
$A_1^*A_1=A_2^*A_2=P^*P=qI$ and
$A_1^*A_2=P^*M_{1|2}P=-\frac{P^*P}{p-1}+\frac{pP^*P\,P}{(p-1)(p-2)}
=-\frac q{p-1}I+\frac P{(p-1)(p-2)}$, since $pq=1$.
For $A_3$ write $\Psi_h(a,b)=\E_s\chi(s-a)\chi(s-b)h(s)$, so $A_3h(x)=-\Psi_h(x_2,x_1)$
on $H_0$, and $\varphi_s(b)=\chi(s-b)$. By \ref{it:paley:corr},
$(P\varphi_s)(a)=q-\one[a=s]$, and $\Jproj\varphi_s=0$, so
\[
 \E[\chi(s-x_1)\mid x_2=a]=(M_{1|2}\varphi_s)(a)
 =-\frac{\chi(s-a)}{p-1}-\frac{p\one[a=s]-1}{(p-1)(p-2)} .
\]
Multiplying by $\chi(s-a)h(s)$ and averaging over $s$, with $\chi(s-a)^2=\one[s\ne a]$
and $\chi(s-a)\one[a=s]=0$, gives
$\E[\Psi_h(a,x_1)\mid x_2=a]=\frac q{p-1}h(a)+\frac{Ph(a)}{(p-1)(p-2)}$ for $h\in H_0$,
which is $-V_2^*A_3h(a)$. Conditioning on $x_1$ with $M_{1|2}^*$ gives $V_1^*A_3$ in
the same way. Then $A_1^*A_3=P^*V_2^*A_3$ and $A_2^*A_3=P^*V_1^*A_3$, and $P^*=-P$,
$P^2=-qI$ give the stated values. For $A_3^*A_3$: the law of $(x_2,x_1)$ is
$(p-2-\chi(a-b))/(p(p-1)(p-2))$ on distinct pairs, $\Psi_h$ is symmetric and the
antisymmetric part $\chi(a-b)$ integrates to zero against
$\Psi_{h'}\Psi_h$; with $\E_{a,b}\Psi_{h'}\Psi_h=\E_{s,s'}h'(s)h(s')(\one[s=s']-q)^2
=q(1-2q)\ip{h'}h$ and $\Psi_h(a,a)=-qh(a)$, one gets
$\frac p{p-1}\bigl[q(1-2q)-q^3\bigr]\ip{h'}h=\frac{1-2q-q^2}{p-1}\ip{h'}h$, which is at
most $q\ip hh$ for $h'=h$. Next $V_0^*A_1=V_0^*V_2P=(P^2-qP)/(1-q)=-(I+P)/(p-1)$
and $V_0^*A_2=V_0^*V_1P=-\frac P{p-1}-\frac{pP^2}{(p-1)(p-2)}$. Finally, given
$x_0=s$ and $x_2=a$, $\E[\chi(t-x_1)\mid x_0,x_2]=-(\chi(t-s)+\chi(t-a))/(p-2)$; so
\begin{align*}
 \E\bigl[\chi(t-x_2)\chi(t-x_1)\bigm|x_0=s\bigr]
 &=-\frac{\chi(t-s)\E[\chi(t-x_2)\mid x_0=s]+\Pr(x_2\ne t\mid x_0=s)}{p-2}\\
 &=-\frac{1-2\cdot\one[t\ne s]/(p-1)}{p-2},
\end{align*}
by $\E[\chi(t-x_2)\mid x_0=s]=(1-\chi(t-s)-p\one[t=s])/(p-1)$ (the computation of
$V_0^*V_2$ applied to $\varphi_t$) and
$\Pr(x_2=t\mid x_0=s)=(\one[t\ne s]+\chi(t-s))/(p-1)$. Averaging against $h\in H_0$
gives $V_0^*A_3h=-\E[\Psi_h(x_2,x_1)\mid x_0]=\frac{2q}{(p-1)(p-2)}h$.

The formula for $V_0^*A_\xi$ is linearity with $1/(p-1)=q/(1-q)$ and
$1/((p-1)(p-2))=q^2/((1-q)(1-2q))$. For the norm, the Schur test applied to the
matrix of the norms $\|A_i^*A_j\|$ (with $\|P\|=r$) bounds
$\|(A_1,A_2,A_3)\|^2$ by its largest row sum. The first two rows sum to
$q+\frac{q^2}{1-q}+\frac{rq^2}{(1-q)(1-2q)}+\frac{q^3}{(1-q)(1-2q)}+\frac{q^2r}{1-q}
=q\bigl(1+\frac{q(1+2r)}{1-2q}\bigr)=q\gamma^2$, and the third row is smaller by
$\frac{2q^2(1-2q+qr-q^2)}{(1-q)(1-2q)}>0$.

\ref{it:paley:perturb} $|\ip FG_m-\ip FG_{\nu_F}|=|\E_{\nu_F}[(g-1)\overline FG]|\le
e\,\E_{\nu_F}|F||G|$, and Cauchy--Schwarz gives the first claim; the second is
the case $F=G$. The coordinate marginals of $m$ are uniform, so
$\|V_if\|_m=\|f\|=\|V_if\|_{\nu_F}$ and the means of $V_if$, $f\in H_0$, vanish under
both laws. For the adjoints,
$\ip f{(V_{i,m}^*-V_{i,\nu_F}^*)Ah}=\ip{V_if}{Ah}_m-\ip{V_if}{Ah}_{\nu_F}$, which is at
most $e\|f\|\|A\|\|h\|$ in modulus.
\end{proof}

In the middle range these identities enter three times: the counts of
\ref{it:paley:counts} make the middle pair law exact
(\cref{lem:record:class-three}) and give the actual pair laws of the fold
records (\cref{lem:actual-features}); the Grams of \ref{it:paley:triples} give
the channel norms of the reference fold record (\cref{lem:old-f-law}) and, with
\ref{it:paley:perturb}, the block-extension estimates for the actual features (\cref{lem:oct24:feature-block});
and \ref{it:paley:features} is the Gram-matrix calculus behind the
joint bound for the reference and the signed comparison fold
records (\cref{lem:record:direct-gram}), with $A_3$ the second-order feature of the
comparison record.

The matching records of the middle range are uniform on pairs of disjoint
injective triples, and at the reserved cut their Markov operator relative to
the actual half law is the following operator (the identification is \cref{lem:record:gram}). Its negativity is exact and
uses no character sum.

\begin{lemma}[The disjoint-triple operator]\label{sec:frontier:empty-triple-negative}
Let $p\ge7$, let $\mathcal T$ be the set of ordered triples of distinct elements
of $\Fp$ with its uniform law, and let $T$ be the Markov operator on
$L^2(\mathcal T)$ which, given a triple $x$, chooses a triple $y\in\mathcal T$
disjoint from $x$ uniformly:
$Tf(x)=\frac1{(p-3)(p-4)(p-5)}\sum_{y\cap x=\emptyset}f(y)$.
\begin{enumerate}[label=\textup{(\alph*)}]
\item\label{it:triple:spec} $T$ is self-adjoint, annihilates the functions whose
 average over the orderings of each unordered triple vanishes, and on functions
 of the unordered triple it has the eigenvalues
 \begin{equation}\label{eq:triple:spectrum}
  \lambda_j=(-1)^j\binom{p-3-j}{3-j}\Big/\binom{p-3}3\qquad(j=0,1,2,3),
 \end{equation}
 that is $1$, $-\frac3{p-3}$, $\frac6{(p-3)(p-4)}$ and
 $-\frac6{(p-3)(p-4)(p-5)}$, the eigenvalue $1$ having the constants as its
 eigenspace.
\item\label{it:triple:bound} Hence $-\frac3{p-3}I\preceq T\preceq I$, the lower constant
 is attained, and $\|T-\Jproj\|=\frac3{p-3}$.
\item\label{it:triple:block} For a block $\Blk$ of primes with least prime $b\ge7$, the
 tensor product $T_{\Blk}=\bigotimes_{p\in\Blk}T$ on $L^2$ of the product of the
 uniform laws satisfies $-\frac3{b-3}I\preceq T_{\Blk}\preceq I$ and
 $\|T_{\Blk}-\Jproj\|\le\frac3{b-3}$.
\end{enumerate}
\end{lemma}

\begin{proof}
\ref{it:triple:spec} The kernel of $T$ with respect to the uniform law is symmetric,
so $T$ is self-adjoint, and $T\one=\one$. For $f$ on $\mathcal T$ let $\bar f$ be its
average over the six orderings of each triple. Then
$Tf(x)=\binom{p-3}3^{-1}\sum_Y\bar f(Y)$, the sum over the unordered triples $Y$
disjoint from $\{x_0,x_1,x_2\}$; so $Tf=0$ if $\bar f=0$, and on the functions of
the unordered triple $T$ acts as the operator $\bar T$ on functions of
$3$-subsets of $\Fp$ given by the same formula.

For $0\le j\le3$ let $\mathcal F_j$ be the space of functions on $j$-subsets, with
the counting inner product. Put $(U_jf)(B)=\sum_{x\in B}f(B\setminus\{x\})$ for
$f\in\mathcal F_j$ and $|B|=j+1$, and $D_{j+1}=U_j^*$, so
$(D_{j+1}g)(A)=\sum_{y\notin A}g(A\cup\{y\})$. Counting the terms that add and
remove the same element gives $D_{j+1}U_j-U_{j-1}D_j=(p-2j)I$ on $\mathcal F_j$
($U_{-1}D_0=0$). For $j\le2<p/2$ this gives $D_{j+1}U_j\succeq(p-2j)I$, so $U_j$ is
injective and $\mathcal F_{j+1}=U_j(\mathcal F_j)\oplus\ker D_{j+1}$ (the range of
$U_j$ is orthogonal to $\ker U_j^*$). Iterating, $\mathcal F_3$ is the direct sum
of the four spaces of \emph{lifts} $f_h(Y)=\sum_{A\subseteq Y,|A|=j}h(A)$ of
\emph{harmonic} functions $h\in\ker D_j$ ($j=0,\dots,3$, with $\ker D_0=\mathcal F_0$
the constants); indeed the lift $U_2\cdots U_jh$ is $(3-j)!\,f_h$.

If $h$ is harmonic on $j$-subsets, the sum of $h(A)$ over the $j$-sets $A$
containing a fixed set $B$ with $|B|<j$ is zero: for $|B|=j-1$ it is $(D_jh)(B)$,
and for smaller $B$ it is $(j-|B|)^{-1}$ times the sum of these over the
$(j-1)$-sets $B'\supseteq B$, since each $A\supseteq B$ contains $j-|B|$ of them.
Hence inclusion--exclusion over the elements of a $3$-set $S$ gives
\[
 \sum_{|A|=j,\ A\cap S=\emptyset}h(A)=\sum_{B\subseteq S}(-1)^{|B|}\sum_{A\supseteq B}h(A)
 =(-1)^j\sum_{A\subseteq S,\ |A|=j}h(A) .
\]
Each $j$-set $A$ disjoint from $S$ lies in $\binom{p-3-j}{3-j}$ of the $3$-sets
disjoint from $S$, so
$\sum_{Y\cap S=\emptyset}f_h(Y)=\binom{p-3-j}{3-j}(-1)^jf_h(S)$, that is
$\bar Tf_h=\lambda_jf_h$. The values \eqref{eq:triple:spectrum} follow by
evaluating the binomial coefficients; the four eigenvalues are distinct, so the
four spaces are the eigenspaces of the self-adjoint operator $\bar T$, and the
$j=0$ space is the constants.

\ref{it:triple:bound} The spectrum of $T$ lies in $\{1,\lambda_1,\lambda_2,\lambda_3,0\}$. For
$p\ge7$, $|\lambda_3|\le\lambda_2\le\frac3{p-3}=|\lambda_1|$ (these are $1\le p-5$ and
$2\le p-4$), and the $j=1$ space is not zero.

\ref{it:triple:block} The eigenvalues of $T_{\Blk}$ are the products of eigenvalues of
the factors, with tensor products of eigenvectors. Every factor has modulus at
most one. A negative product has a negative factor, of modulus at most
$\frac3{p-3}\le\frac3{b-3}$; and a product other than the one with all factors
equal to $1$ (whose eigenvector is the constant) has a factor of modulus at most
$\frac3{p-3}$.
\end{proof}

The operator $\bar T$ is the normalized adjacency operator of the Kneser graph
of $3$-subsets of a $p$-set, two subsets being adjacent when they are disjoint,
and \eqref{eq:triple:spectrum} is its classical spectrum; the proof above is
self-contained. In the middle range (\cref{lem:record:gram}) the bound is used as follows: for a block with least prime $b\ge7$
(writing $r=b^{-1/2}$ in this paragraph), an unfolded
empty-triple record of mass $w_Er$ has reserved
negative part at least $-\frac3{b-3}$ times its own half diagonal, so it is
compensated by a folded actual-half diagonal of mass
$\frac3{b-3}w_Er=\frac{3w_E}{1-3/b}r^3$; in particular the doubled correction
used to replace a removed folded empty-triple diagonal has mass at most
$\frac{6w_E}{1-3/b}r^3$. The actual half law of the record is exactly the
uniform law on injective triples, so no change of reference density is
required.

\paragraph{Consequences and what to read next.}
Weil's bound in odd and even degree (\cref{thm:weil}) and the two-root identity
\eqref{eq:two-root}; the square shift and the Paley operator
(\cref{lem:square-shift}); square-affine invariance and its consequences for one-slot
and pair marginals (\cref{lem:marginals}); the conditional operator of
\cref{def:mixing}; and the admissible majorants (\cref{lem:majorant}) for the
transfer of thresholds.
